\chapter{Chapter03 — Program}
%%% generated by the blueprint pipeline from FABL.Chapter03.LearningTheory.Program
\section{Overview}

Declarations of the Lean module \texttt{FABL.Chapter03.LearningTheory.Program}.

\section{Declarations}

\begin{theorem}[Uniform distribution under bijection]
\lean{FABL.map_uniformPMF_equiv}
\label{FABL.map_uniformPMF_equiv}
\leanok
\uses{FABL.uniformPMF}
\difficulty{4}
Let $\\alpha$ and $\\beta$ be finite nonempty sets, and let $e:\\alpha\\to\\beta$ be a
bijection. Mapping a uniformly chosen element of $\\alpha$ through $e$ gives the uniform
distribution on $\\beta$.
\end{theorem}

\begin{definition}[Learning accuracy range]
\lean{FABL.LearningAccuracy}
\statementsource{Analysis-of-Boolean-Functions-by-Ryan-ODonnell}{
  pdf-ffbe9c7b3306-p079-b001,
  pdf-ffbe9c7b3306-p089-b011,
  pdf-ffbe9c7b3306-p090-b001
}
\label{FABL.LearningAccuracy}
\leanok
Learning accuracy is the set of rational numbers $[0,1/2]$.
\end{definition}

\begin{definition}[Positive learning parameter range]
\lean{FABL.PositiveLearningParameter}
\statementsource{Analysis-of-Boolean-Functions-by-Ryan-ODonnell}{pdf-ffbe9c7b3306-p079-b001}
\label{FABL.PositiveLearningParameter}
\leanok
A positive learning parameter is a rational number strictly greater than $0$ and at most
$1/2$.
\end{definition}

\begin{theorem}[Real-valued learning parameter bounds]
\lean{FABL.positiveLearningParameter_toReal_mem_Ioc}
\label{FABL.positiveLearningParameter_toReal_mem_Ioc}
\leanok
\uses{FABL.PositiveLearningParameter}
\difficulty{3}
If a rational learning parameter $\\varepsilon$ satisfies $0<\\varepsilon\\leq
\\tfrac12$, then its image in $\\mathbb R$ also lies in $(0,\\tfrac12]$.
\end{theorem}

\begin{theorem}[Rational learning parameter approximation]
\lean{FABL.exists_positiveLearningParameter_between_half}
\label{FABL.exists_positiveLearningParameter_between_half}
\leanok
\uses{FABL.PositiveLearningParameter}
\difficulty{4}
For every real number $x>0$, there is a rational number $q$ with $0<q\leq \tfrac12$ such
that $\tfrac12\min\{x,\tfrac12\}<q<\min\{x,\tfrac12\}$.
\end{theorem}

\begin{definition}[Fourier estimator failure bits]
\lean{FABL.fourierEstimatorFailureBits}
\statementsource{Analysis-of-Boolean-Functions-by-Ryan-ODonnell}{pdf-ffbe9c7b3306-p080-b002}
\label{FABL.fourierEstimatorFailureBits}
\leanok
\uses{FABL.PositiveLearningParameter}
For a positive learning parameter $$, define the failure-bit count to be the base-$2$
ceiling logarithm of the ceiling of $2/$.
\end{definition}

\begin{definition}[Fourier estimator sample count]
\lean{FABL.fourierEstimatorSampleCount}
\statementsource{Analysis-of-Boolean-Functions-by-Ryan-ODonnell}{
  pdf-ffbe9c7b3306-p080-b002,
  pdf-ffbe9c7b3306-p080-b003
}
\label{FABL.fourierEstimatorSampleCount}
\leanok
\uses{FABL.PositiveLearningParameter, FABL.fourierEstimatorFailureBits}
For rational parameters $\varepsilon,\delta$ with $0<\varepsilon,\delta\leq 1/2$, the
Fourier estimator sample count is the least natural number at least
$2B(\delta)/\varepsilon^2$, where $B(\delta)$ denotes the quantity assigned to $\delta$
by the estimator’s failure-bit function.
\end{definition}

\begin{theorem}[Positivity of the failure-bit count]
\lean{FABL.fourierEstimatorFailureBits_pos}
\label{FABL.fourierEstimatorFailureBits_pos}
\leanok
\uses{FABL.PositiveLearningParameter, FABL.fourierEstimatorFailureBits}
\difficulty{4}
Let $n$ be a nonnegative integer and let $\delta\in\mathbb Q$ satisfy $0<\delta\le 1/2$.
The integer $b(\delta)=\lceil\log_2\lceil 2/\delta\rceil\rceil$ is positive.
\end{theorem}

\begin{theorem}[Positivity of the sample count]
\lean{FABL.fourierEstimatorSampleCount_pos}
\label{FABL.fourierEstimatorSampleCount_pos}
\leanok
\uses{FABL.PositiveLearningParameter, FABL.fourierEstimatorFailureBits, FABL.fourierEstimatorFailureBits_pos, FABL.fourierEstimatorSampleCount}
\difficulty{3}
Let $n$ be a nonnegative integer and let $\varepsilon,\delta\in\mathbb Q$ satisfy
$0<\varepsilon,\delta\le 1/2$. With $b(\delta)=\lceil\log_2\lceil 2/\delta\rceil\rceil$
and $m=\lceil 2b(\delta)/\varepsilon^2\rceil$, one has $m>0$.
\end{theorem}

\begin{theorem}[Failure-bit probability bound]
\lean{FABL.two_div_pow_fourierEstimatorFailureBits_le}
\label{FABL.two_div_pow_fourierEstimatorFailureBits_le}
\leanok
\uses{FABL.PositiveLearningParameter, FABL.fourierEstimatorFailureBits}
\difficulty{4}
Let $\delta$ be a rational number with $0<\delta\le 1/2$, and let $k=\lceil\log_2\lceil
2/\delta\rceil\rceil$. Then $2/2^k\le\delta$.
\end{theorem}

\begin{theorem}[Failure bits bounded by the sample exponent]
\lean{FABL.fourierEstimatorFailureBits_le_sampleExponent}
\label{FABL.fourierEstimatorFailureBits_le_sampleExponent}
\leanok
\uses{FABL.PositiveLearningParameter, FABL.fourierEstimatorFailureBits, FABL.fourierEstimatorSampleCount}
\difficulty{4}
Let $\varepsilon,\delta$ be rational parameters in $(0,1/2]$. If the failure-bit count
is $b=\lceil\log_2\lceil 2/\delta\rceil\rceil$ and the sample count is $N=\lceil
2b/\varepsilon^2\rceil$, then $b\leq N\varepsilon^2/2$.
\end{theorem}

\begin{theorem}[Upper bound for the sample count]
\lean{FABL.fourierEstimatorSampleCount_cast_le}
\label{FABL.fourierEstimatorSampleCount_cast_le}
\leanok
\uses{FABL.PositiveLearningParameter, FABL.fourierEstimatorFailureBits, FABL.fourierEstimatorFailureBits_pos, FABL.fourierEstimatorSampleCount}
\difficulty{4}
Let $n$ be a nonnegative integer and let $\varepsilon,\delta\in\mathbb Q$ satisfy
$0<\varepsilon,\delta\le 1/2$. Define $b(\delta)=\lceil\log_2\lceil
2/\delta\rceil\rceil$ and $m=\lceil 2b(\delta)/\varepsilon^2\rceil$. Then $m\le
4b(\delta)/\varepsilon^2$.
\end{theorem}

\begin{theorem}[Fourier estimator failure bound]
\lean{FABL.two_mul_exp_neg_fourierEstimatorSampleCount_le}
\label{FABL.two_mul_exp_neg_fourierEstimatorSampleCount_le}
\leanok
\uses{FABL.PositiveLearningParameter, FABL.fourierEstimatorFailureBits, FABL.fourierEstimatorFailureBits_le_sampleExponent, FABL.fourierEstimatorSampleCount, FABL.two_div_pow_fourierEstimatorFailureBits_le}
\difficulty{5}
Let $\bbr$-valued learning parameters $bfepsilon,bfd$ lie in $(0,1/2]$. Let $B$ be the
least integer $m$ such that $2^ma0\text{?}$
\end{theorem}

\begin{theorem}[Real interval for learning accuracy]
\lean{FABL.learningAccuracy_toReal_mem_Icc}
\label{FABL.learningAccuracy_toReal_mem_Icc}
\leanok
\uses{FABL.LearningAccuracy}
\difficulty{3}
If $\varepsilon$ is a rational number in the interval $[0,1/2]$, then its image in
$\mathbb R$ also belongs to $[0,1/2]$.
\end{theorem}

\begin{definition}[Learning access modes]
\lean{FABL.LearningAccess}
\statementsource{Analysis-of-Boolean-Functions-by-Ryan-ODonnell}{
  pdf-ffbe9c7b3306-p079-b001,
  pdf-ffbe9c7b3306-p088-b004,
  pdf-ffbe9c7b3306-p089-b004
}
\label{FABL.LearningAccess}
\leanok
A learning procedure may receive independent uniformly distributed input-label pairs, or
it may request the target’s value at an input of its choice.
\end{definition}

\begin{definition}[Learning cost]
\lean{FABL.LearningCost}
\statementsource{Analysis-of-Boolean-Functions-by-Ryan-ODonnell}{
  pdf-ffbe9c7b3306-p079-b001,
  pdf-ffbe9c7b3306-p079-b002
}
\label{FABL.LearningCost}
\leanok
A learning cost is a triple of natural numbers recording, respectively, the number of
random-example oracle calls, the number of membership queries, and the number of
explicitly charged local computation steps.
\end{definition}

\begin{definition}[Learning cost as a triple]
\lean{FABL.LearningCost.toTriple}
\label{FABL.LearningCost.toTriple}
\leanok
\uses{FABL.LearningCost}
A learning cost is represented by the ordered triple whose entries are, respectively,
its number of random-example oracle calls, its number of membership queries, and its
number of explicitly charged local computation steps.
\end{definition}

\begin{theorem}[Injectivity of the cost triple]
\lean{FABL.LearningCost.toTriple_injective}
\label{FABL.LearningCost.toTriple_injective}
\leanok
\uses{FABL.LearningCost.toTriple}
\difficulty{4}
The map sending a learning cost to the ordered triple consisting of its number of
random-example calls, membership queries, and charged local computation steps is
injective.
\end{theorem}

\begin{definition}[Cost of a random bit]
\lean{FABL.LearningCost.coin}
\label{FABL.LearningCost.coin}
\leanok
\uses{FABL.LearningCost}
The cost of drawing one unbiased random bit is zero random-example calls, zero
membership queries, and one unit of local computation.
\end{definition}

\begin{definition}[Cost of a random example]
\lean{FABL.LearningCost.example}
\statementsource{Analysis-of-Boolean-Functions-by-Ryan-ODonnell}{
  pdf-ffbe9c7b3306-p079-b001,
  pdf-ffbe9c7b3306-p089-b004,
  pdf-ffbe9c7b3306-p155-b010
}
\label{FABL.LearningCost.example}
\leanok
\uses{FABL.LearningCost}
The cost assigned to one random example consists of one random-example oracle call, no
membership queries, and one explicitly charged local computation step.
\end{definition}

\begin{definition}[Cost of a membership query]
\lean{FABL.LearningCost.query}
\statementsource{Analysis-of-Boolean-Functions-by-Ryan-ODonnell}{
  pdf-ffbe9c7b3306-p079-b001,
  pdf-ffbe9c7b3306-p088-b004,
  pdf-ffbe9c7b3306-p089-b004
}
\label{FABL.LearningCost.query}
\leanok
\uses{FABL.LearningCost}
The cost assigned to one membership query consists of no random-example oracle calls,
one membership query, and one explicitly charged local computation step.
\end{definition}

\begin{definition}[Learning programs]
\lean{FABL.LearningProgram}
\label{FABL.LearningProgram}
\leanok
\uses{FABL.LearningAccess, FABL.Sign}
For a dimension $n\in\mathbb{N}$, an access mode (random-example access or
membership-query access), and an output type, a learning program may return a value,
draw an unbiased random bit, draw one or a batch of independent uniform labeled
examples, draw a batch of uniform cube inputs, or draw a finite matrix of independent
uniform labeled examples. It may also compute a pure output from such a matrix with an
explicitly specified local-work charge, query the target at a chosen cube input, or
charge a specified number of local computation steps; the random-example and query
operations are available only with their respective access modes.
\end{definition}

\begin{definition}[Add cost to an outcome]
\lean{FABL.LearningProgram.addOutcomeCost}
\label{FABL.LearningProgram.addOutcomeCost}
\leanok
\uses{FABL.LearningCost}
Given an output type, an extra learning cost, and an outcome consisting of an output and
a learning cost, adding cost to the outcome preserves the output and adds the extra cost
to its existing cost.
\end{definition}

\begin{definition}[Bind learning programs]
\lean{FABL.LearningProgram.bind}
\label{FABL.LearningProgram.bind}
\leanok
\uses{FABL.LearningAccess, FABL.LearningProgram}
Fix a dimension, an access mode, and input and output types, and let a continuation
assign an output-type program to each input value. Binding the continuation to a program
replaces its returned value with the corresponding continuation program, propagates
through random-bit draws, example draws, input batches, queries, and local-work charges,
and preserves those operations. For a program that computes an output from a
random-example matrix with a specified work charge, binding draws the same matrix,
charges that work, and then runs the continuation on the computed output.
\end{definition}

\begin{definition}[Map a function over a program]
\lean{FABL.LearningProgram.map}
\label{FABL.LearningProgram.map}
\leanok
\uses{FABL.LearningAccess, FABL.LearningProgram, FABL.LearningProgram.bind}
Fix a dimension, an access mode, input and output types, a function from the input type
to the output type, and a program returning an input. Mapping the function over the
program returns its image under the function while retaining the program's operations
and costs.
\end{definition}

\begin{definition}[Sequence learning programs]
\lean{FABL.LearningProgram.sequence}
\label{FABL.LearningProgram.sequence}
\leanok
\uses{FABL.LearningAccess, FABL.LearningProgram, FABL.LearningProgram.bind, FABL.LearningProgram.map}
For a dimension, an access mode, and an output type, sequencing a list of programs that
each return that type produces a program returning the list of their outputs in order.
\end{definition}

\begin{definition}[Random example matrix outcome]
\lean{FABL.LearningProgram.randomExampleMatrixOutcome}
\statementsource{Analysis-of-Boolean-Functions-by-Ryan-ODonnell}{
  pdf-ffbe9c7b3306-p078-b006,
  pdf-ffbe9c7b3306-p079-b001
}
\label{FABL.LearningProgram.randomExampleMatrixOutcome}
\leanok
\uses{FABL.BooleanFunction, FABL.LearningCost, FABL.Sign}
Let $X=\{-1,1\}^n$, let $f:X\to\{-1,1\}$ be a target function, let $I$ be a finite set,
and let $m,w$ be natural numbers. Given an array of inputs $x_{k,i}\in X$ indexed by
$k\in I$ and $1\leq i\leq m$, apply a specified function to the array of labeled pairs
$(x_{k,i},f(x_{k,i}))$. The resulting outcome is that function’s value together with the
cost $(|I|m,0,|I|m+w)$, recording random examples, membership queries, and local
computation steps in that order.
\end{definition}

\begin{definition}[Uniform input matrix law]
\lean{FABL.LearningProgram.randomExampleMatrixInputLaw}
\statementsource{Analysis-of-Boolean-Functions-by-Ryan-ODonnell}{
  pdf-ffbe9c7b3306-p078-b006,
  pdf-ffbe9c7b3306-p079-b001
}
\label{FABL.LearningProgram.randomExampleMatrixInputLaw}
\leanok
\uses{FABL.uniformPMF}
Let $I$ be a finite index set and let $m,n$ be natural numbers. The random-example input
matrix law is the uniform probability mass function on arrays of points of $\{-1,1\}^n$,
indexed by $k\in I$ and $1\leq i\leq m$.
\end{definition}

\begin{theorem}[Uniform law of matrix inputs]
\lean{FABL.LearningProgram.randomExampleMatrixInputLaw_eq_uniformPMF}
\label{FABL.LearningProgram.randomExampleMatrixInputLaw_eq_uniformPMF}
\leanok
\uses{FABL.LearningProgram.randomExampleMatrixInputLaw, FABL.uniformPMF}
\difficulty{3}
Let $n,m$ be natural numbers, let $\kappa$ be a type with decidable equality, and let
$I$ be a finite subset of $\kappa$. The probability mass function for the matrix of
inputs indexed by $I$ and $\mathrm{Fin}(m)$ is the uniform probability mass function on
all such matrices of $n$-dimensional sign vectors.
\end{theorem}

\begin{definition}[Random example matrix output law]
\lean{FABL.LearningProgram.randomExampleMatrixOutputLaw}
\statementsource{Analysis-of-Boolean-Functions-by-Ryan-ODonnell}{
  pdf-ffbe9c7b3306-p078-b006,
  pdf-ffbe9c7b3306-p079-b001
}
\label{FABL.LearningProgram.randomExampleMatrixOutputLaw}
\leanok
\uses{FABL.BooleanFunction, FABL.LearningCost, FABL.LearningProgram.randomExampleMatrixInputLaw, FABL.LearningProgram.randomExampleMatrixOutcome, FABL.Sign}
Let $X=\{-1,1\}^n$, let $f:X\to\{-1,1\}$ be a target function, let $I$ be a finite index
set, and let $m,w$ be natural numbers. Choose an array of inputs indexed by $I$ and
$\{1,\ldots,m\}$ uniformly from all such arrays, label each input $x$ by $f(x)$, and
apply a specified function to the resulting array. The output law is the distribution of
the pair consisting of that function’s value and the cost $(|I|m,0,|I|m+w)$, recording
random examples, membership queries, and local computation steps in that order.
\end{definition}

\begin{theorem}[Output law on events]
\lean{FABL.LearningProgram.randomExampleMatrixOutputLaw_toOuterMeasure_apply}
\label{FABL.LearningProgram.randomExampleMatrixOutputLaw_toOuterMeasure_apply}
\leanok
\uses{FABL.BooleanFunction, FABL.LearningCost, FABL.LearningProgram.randomExampleMatrixInputLaw, FABL.LearningProgram.randomExampleMatrixOutcome, FABL.LearningProgram.randomExampleMatrixOutputLaw, FABL.Sign}
\difficulty{3}
Let $n,m,w$ be natural numbers, let $\kappa$ be a type, let $I$ be a finite subset of
$\kappa$, and let $f$ be a Boolean function on the $n$-dimensional sign cube. For any
function $g$ from matrices of inputs and their labels, indexed by $I$ and
$\mathrm{Fin}(m)$, to a type $\alpha$, and any event $E$ in $\alpha$ paired with a
learning cost, the output law assigns $E$ the same outer measure as the input law
assigns the set of input matrices whose labeled output, with its specified learning
cost, lies in $E$.
\end{theorem}

\begin{definition}[Run a program with cost]
\lean{FABL.LearningProgram.runWithCost}
\label{FABL.LearningProgram.runWithCost}
\leanok
\uses{FABL.BooleanFunction, FABL.LearningAccess, FABL.LearningCost, FABL.LearningCost.coin, FABL.LearningCost.example, FABL.LearningCost.query, FABL.LearningProgram, FABL.LearningProgram.addOutcomeCost, FABL.LearningProgram.randomExampleMatrixOutputLaw, FABL.uniformPMF}
Fix a dimension, an access mode, an output type, and a Boolean target function on the
$n$-dimensional sign cube. Running a learning program with cost produces a probability
distribution on pairs of outputs and learning costs: a pure return has zero cost; a
random bit costs one unit of local work; each random labeled example costs one example
call and one unit of work; a batch of $m$ labeled examples costs $m$ example calls and
$m$ units of work; a batch of $m$ uniform cube inputs costs $mn$ units of work; and a
matrix indexed by a finite set $I$ with $m$ examples per index costs $|I|m$ example
calls and $|I|m$ units of work. A matrix-output operation uses that same matrix cost
plus its specified work charge, a membership query costs one query and one unit of work,
and an explicitly charged computation costs its specified number of work units.
\end{definition}

\begin{theorem}[Zero added outcome cost]
\lean{FABL.LearningProgram.addOutcomeCost_zero}
\label{FABL.LearningProgram.addOutcomeCost_zero}
\leanok
\uses{FABL.LearningCost, FABL.LearningProgram.addOutcomeCost}
\difficulty{4}
For any type of outcomes and any outcome paired with a learning cost, adding zero to
that cost leaves the pair unchanged.
\end{theorem}

\begin{theorem}[Zero added cost preserves the distribution]
\lean{FABL.LearningProgram.map_addOutcomeCost_zero}
\label{FABL.LearningProgram.map_addOutcomeCost_zero}
\leanok
\uses{FABL.LearningCost, FABL.LearningProgram.addOutcomeCost, FABL.LearningProgram.addOutcomeCost_zero}
\difficulty{3}
For any type of outputs and any probability mass function on pairs consisting of an
output and a learning cost, mapping each pair to the same output with zero added to its
cost leaves the probability mass function unchanged.
\end{theorem}

\begin{theorem}[Commuting cost addition with bind and map]
\lean{FABL.LearningProgram.bind_map_addOutcomeCost_comm}
\label{FABL.LearningProgram.bind_map_addOutcomeCost_comm}
\leanok
\uses{FABL.LearningCost, FABL.LearningProgram.addOutcomeCost}
\difficulty{4}
Let $p$ be a probability mass function on pairs consisting of an outcome in $\alpha$ and
a learning cost, let $q$ assign to each outcome in $\alpha$ a probability mass function
on pairs consisting of an outcome in $\beta$ and a learning cost, and fix an additional
learning cost $e$. Binding $p$ to the law obtained by adding the first-stage cost to
each cost from $q$, then adding $e$, gives the same law as first adding $e$ to the cost
drawn from $p$ and then binding to the law obtained by adding that first-stage cost to
each cost from $q$.
\end{theorem}

\begin{theorem}[Cost law for program sequencing]
\lean{FABL.LearningProgram.runWithCost_bind}
\label{FABL.LearningProgram.runWithCost_bind}
\leanok
\uses{FABL.BooleanFunction, FABL.LearningAccess, FABL.LearningCost.coin, FABL.LearningCost.example, FABL.LearningCost.query, FABL.LearningProgram, FABL.LearningProgram.addOutcomeCost, FABL.LearningProgram.bind, FABL.LearningProgram.bind_map_addOutcomeCost_comm, FABL.LearningProgram.map_addOutcomeCost_zero, FABL.LearningProgram.randomExampleMatrixInputLaw, FABL.LearningProgram.randomExampleMatrixOutcome, FABL.LearningProgram.randomExampleMatrixOutputLaw, FABL.LearningProgram.runWithCost}
\difficulty{5}
Let $n$ be a natural number, fix a learning-access mode, and let $f$ assign to each
output $a$ a learning program with output in $\beta$. For any Boolean target function on
the $n$-dimensional sign cube and any program with output in $\alpha$, the law of
running the program obtained by binding $f$ to it equals the law obtained by first
running the original program and then, for each outcome, running $f$ on its output and
adding the original outcome's learning cost to the subsequent cost.
\end{theorem}

\begin{theorem}[Cost law under output mapping]
\lean{FABL.LearningProgram.runWithCost_map}
\label{FABL.LearningProgram.runWithCost_map}
\leanok
\uses{FABL.BooleanFunction, FABL.LearningAccess, FABL.LearningProgram, FABL.LearningProgram.map, FABL.LearningProgram.runWithCost, FABL.LearningProgram.runWithCost_bind}
\difficulty{3}
Let $n$ be a natural number, fix a learning-access mode, and let
$f\colon\alpha\to\beta$. For any Boolean target function on the $n$-dimensional sign
cube and any program with output in $\alpha$, running the program with its output mapped
by $f$ has the same law as mapping each run's output by $f$ while leaving its learning
cost unchanged.
\end{theorem}

\begin{theorem}[Cost law for a sequence of programs]
\lean{FABL.LearningProgram.runWithCost_sequence_cons}
\label{FABL.LearningProgram.runWithCost_sequence_cons}
\leanok
\uses{FABL.BooleanFunction, FABL.LearningAccess, FABL.LearningProgram, FABL.LearningProgram.runWithCost, FABL.LearningProgram.runWithCost_bind, FABL.LearningProgram.runWithCost_map, FABL.LearningProgram.sequence}
\difficulty{3}
Let $n$ be a natural number, fix a learning-access mode, and let $P$ be a program with
output in $\alpha$, followed by a finite list of such programs. For any Boolean target
function on the $n$-dimensional sign cube, the law of running the sequence beginning
with $P$ is obtained by first running $P$, then running the sequence of remaining
programs, returning the list formed by placing the first output before the remaining
outputs, and adding the two learning costs.
\end{theorem}

\begin{theorem}[Cost bound for a sequence of programs]
\lean{FABL.LearningProgram.costProjection_le_of_mem_support_sequence}
\label{FABL.LearningProgram.costProjection_le_of_mem_support_sequence}
\leanok
\uses{FABL.BooleanFunction, FABL.LearningAccess, FABL.LearningCost, FABL.LearningProgram, FABL.LearningProgram.runWithCost, FABL.LearningProgram.runWithCost_sequence_cons, FABL.LearningProgram.sequence}
\difficulty{4}
Let $n$ be a natural number, choose a learning access mode and an output type, and let
$T$ be a Boolean function on the $n$-dimensional sign cube. Suppose $P$ assigns a
natural number to each learning cost, with $P(0)=0$ and $P(c+d)=P(c)+P(d)$ for all
learning costs $c,d$. If, for every program in a finite list and every outcome-cost pair
in the support of that program's run under $T$, the value of $P$ on its cost is at most
$B$, then every outcome-cost pair in the support of the run of the whole sequence has
cost projection at most the length of the list times $B$.
\end{theorem}

\begin{theorem}[Keys recover sequence items on support]
\lean{FABL.LearningProgram.map_output_eq_of_mem_support_sequence}
\label{FABL.LearningProgram.map_output_eq_of_mem_support_sequence}
\leanok
\uses{FABL.BooleanFunction, FABL.LearningAccess, FABL.LearningProgram, FABL.LearningProgram.runWithCost, FABL.LearningProgram.runWithCost_sequence_cons, FABL.LearningProgram.sequence}
\difficulty{4}
Let a target Boolean function, a family of learning programs indexed by a type, and a
key function from outputs to that index type be given. Suppose that for each index,
every output-and-cost pair in the support of the corresponding program's outcome law has
output whose key equals that index. Then every outcome-and-cost pair in the support of
the law for running the sequence of programs indexed by a finite list has an output list
whose keys, in order, equal that list.
\end{theorem}

\begin{definition}[Probability of a program event]
\lean{FABL.LearningProgram.eventProbability}
\label{FABL.LearningProgram.eventProbability}
\leanok
\uses{FABL.BooleanFunction, FABL.LearningAccess, FABL.LearningCost, FABL.LearningProgram, FABL.LearningProgram.runWithCost}
Fix a dimension, an access mode, an output type, a learning program, a Boolean target
function on the $n$-dimensional sign cube, and an event on output-cost pairs. The event
probability is the real-valued measure of the outcomes satisfying that event under the
program's run distribution with the given target.
\end{definition}

\begin{theorem}[Event probability under a mapped law]
\lean{FABL.LearningProgram.eventProbability_eq_toMeasure_real_of_runWithCost_eq_map}
\label{FABL.LearningProgram.eventProbability_eq_toMeasure_real_of_runWithCost_eq_map}
\leanok
\uses{FABL.BooleanFunction, FABL.LearningAccess, FABL.LearningCost, FABL.LearningProgram, FABL.LearningProgram.eventProbability, FABL.LearningProgram.runWithCost}
\difficulty{3}
Let $n$ be a natural number, choose a learning access mode, and let $P$ be a learning
program with output in $\beta$. Suppose $\alpha$ is a measurable space whose singletons
are measurable, $\mu$ is a probability mass function on $\alpha$, $f:\alpha\to\beta$,
and $c$ is a fixed learning cost. If the law of the output and cost from running $P$
equals the law obtained by sampling $x$ from $\mu$ and returning $(f(x),c)$, then for
every predicate $E$ on $\beta$, the probability that the output of $P$ satisfies $E$
equals the measure under $\mu$ of the set of $x$ for which $f(x)$ satisfies $E$.
\end{theorem}

\begin{theorem}[Bad-event bound for a two-stage sample]
\lean{FABL.PMF.toOuterMeasure_bind_le_add}
\label{FABL.PMF.toOuterMeasure_bind_le_add}
\leanok
\difficulty{5}
Let $p$ be a probability mass function on a type $\alpha$, let $q$ assign a probability
mass function on $\beta$ to each element of $\alpha$, and let $A$ and $B$ be bad-event
sets in the respective spaces. If $p$ assigns outer measure at most $\delta_1$ to $A$,
and for every $a$ outside $A$, $q(a)$ assigns outer measure at most $\delta_2$ to $B$,
then the probability mass function obtained by sampling first from $p$ and then from $q$
assigns outer measure at most $\delta_1+\delta_2$ to $B$.
\end{theorem}

\begin{definition}[Outer probability of an event]
\lean{FABL.LearningProgram.outerEventProbability}
\label{FABL.LearningProgram.outerEventProbability}
\leanok
\uses{FABL.BooleanFunction, FABL.LearningAccess, FABL.LearningCost, FABL.LearningProgram, FABL.LearningProgram.runWithCost}
Given a learning program with either random-example access or query access, a Boolean
target on the $n$-dimensional sign cube, and an event in the product of the program’s
output space with its learning cost, its outer event probability is the outer
measure—under the program’s run distribution with that target—of the set of output-cost
pairs satisfying the event.
\end{definition}

\begin{theorem}[Finite outer probability of an event]
\lean{FABL.LearningProgram.outerEventProbability_ne_top}
\label{FABL.LearningProgram.outerEventProbability_ne_top}
\leanok
\uses{FABL.BooleanFunction, FABL.LearningAccess, FABL.LearningCost, FABL.LearningProgram, FABL.LearningProgram.outerEventProbability, FABL.LearningProgram.runWithCost}
\difficulty{4}
For any learning program, target Boolean function, and event on outcome-and-cost pairs,
the event's outer probability is not infinite.
\end{theorem}

\begin{theorem}[Real value of outer event probability]
\lean{FABL.LearningProgram.outerEventProbability_toReal}
\label{FABL.LearningProgram.outerEventProbability_toReal}
\leanok
\uses{FABL.BooleanFunction, FABL.LearningAccess, FABL.LearningCost, FABL.LearningProgram, FABL.LearningProgram.eventProbability, FABL.LearningProgram.outerEventProbability}
\difficulty{1}
For any learning program, target Boolean function, and event on outcome-and-cost pairs,
the real value of the event's outer probability equals its event probability.
\end{theorem}

\begin{theorem}[Monotonicity of event probability]
\lean{FABL.LearningProgram.eventProbability_mono}
\label{FABL.LearningProgram.eventProbability_mono}
\leanok
\uses{FABL.BooleanFunction, FABL.LearningAccess, FABL.LearningCost, FABL.LearningProgram, FABL.LearningProgram.eventProbability, FABL.LearningProgram.outerEventProbability, FABL.LearningProgram.outerEventProbability_ne_top, FABL.LearningProgram.runWithCost}
\difficulty{3}
For any learning program, target Boolean function, and two events on program outcomes
with costs, if every outcome in the first event also belongs to the second, then the
probability of the first event is at most the probability of the second.
\end{theorem}

\begin{theorem}[Bad-event bound for sequential programs]
\lean{FABL.LearningProgram.outerEventProbability_bind_le_add}
\label{FABL.LearningProgram.outerEventProbability_bind_le_add}
\leanok
\uses{FABL.BooleanFunction, FABL.LearningAccess, FABL.LearningProgram, FABL.LearningProgram.addOutcomeCost, FABL.LearningProgram.bind, FABL.LearningProgram.outerEventProbability, FABL.LearningProgram.runWithCost, FABL.LearningProgram.runWithCost_bind, FABL.PMF.toOuterMeasure_bind_le_add}
\difficulty{4}
Let a target Boolean function, a first learning program, and a continuation program for
each possible first output be given. Let the first-stage bad event be a predicate on
first outputs, the final bad event a predicate on final outputs, and let
$\delta_1,\delta_2$ be nonnegative extended real numbers. If the outer probability that
the first output is bad is at most $\delta_1$, and for every first output that is not
bad the outer probability that the continuation's output is bad is at most $\delta_2$,
then the outer probability that the combined program's final output is bad is at most
$\delta_1+\delta_2$.
\end{theorem}

\begin{theorem}[Failure probability bound for sequential programs]
\lean{FABL.LearningProgram.eventProbability_bind_le_add}
\label{FABL.LearningProgram.eventProbability_bind_le_add}
\leanok
\uses{FABL.BooleanFunction, FABL.LearningAccess, FABL.LearningProgram, FABL.LearningProgram.bind, FABL.LearningProgram.eventProbability, FABL.LearningProgram.outerEventProbability, FABL.LearningProgram.outerEventProbability_bind_le_add, FABL.LearningProgram.outerEventProbability_ne_top, FABL.LearningProgram.outerEventProbability_toReal}
\difficulty{4}
Let $T$ be a Boolean function on the $n$-dimensional sign cube, let $P$ be a learning
program with output in $\alpha$, and let $Q$ assign a learning program with output in
$\beta$ to each element of $\alpha$. Suppose $\delta_1,\delta_2\geq 0$, the probability
that the output of $P$ satisfies a predicate $A$ is at most $\delta_1$, and for every
$a\in\alpha$ not satisfying $A$, the probability that the output of $Q(a)$ satisfies a
predicate $B$ is at most $\delta_2$. Then the probability that the output of the program
obtained by binding $P$ to $Q$ satisfies $B$ is at most $\delta_1+\delta_2$.
\end{theorem}

\begin{theorem}[Prepending a good item preserves bad-event probability]
\lean{FABL.LearningProgram.outerEventProbability_map_cons_exists}
\label{FABL.LearningProgram.outerEventProbability_map_cons_exists}
\leanok
\uses{FABL.BooleanFunction, FABL.LearningAccess, FABL.LearningProgram, FABL.LearningProgram.map, FABL.LearningProgram.outerEventProbability, FABL.LearningProgram.runWithCost_map}
\difficulty{3}
Let a target Boolean function, an output type, a fixed head item, and a learning program
returning a list of such items be given. If the head item does not satisfy a specified
bad predicate, then the outer probability that the mapped program's output list contains
any bad item equals the outer probability that the original program's output list
contains any bad item.
\end{theorem}

\begin{theorem}[Union bound for a sequence of programs]
\lean{FABL.LearningProgram.outerEventProbability_sequence_exists_le_nsmul}
\label{FABL.LearningProgram.outerEventProbability_sequence_exists_le_nsmul}
\leanok
\uses{FABL.BooleanFunction, FABL.LearningAccess, FABL.LearningProgram, FABL.LearningProgram.map, FABL.LearningProgram.outerEventProbability, FABL.LearningProgram.outerEventProbability_bind_le_add, FABL.LearningProgram.outerEventProbability_map_cons_exists, FABL.LearningProgram.runWithCost, FABL.LearningProgram.sequence}
\difficulty{4}
Let a target Boolean function, a bad predicate on outputs, and a nonnegative extended
real number $\delta$ be given. For every finite list of learning programs with the same
output type, if each program's outer probability of producing a bad output is at most
$\delta$, then the outer probability that the sequence of programs produces a list
containing at least one bad output is at most the list length times $\delta$.
\end{theorem}
